% ===================================================================
\section{Discussion}
\label{sec:discussion}

\subsection{Why the Bridge Does Not Improve Loss}

The bridge adds 36 parameters per adapter layer. In a rank-24
RhombiLoRA configuration applied to four target modules across 28
transformer layers, this amounts to $28 \times 4 \times 36 = 4{,}032$
bridge parameters---negligible relative to the adapter's total
parameter count ($28 \times 4 \times 2 \times 24 \times d_{\text{model}}$,
where $d_{\text{model}} = 3{,}584$ for Qwen2.5-7B, yielding roughly
19.3M adapter parameters). The bridge contributes 0.02\% of total
adapter capacity.

At this scale, the optimization landscape does not benefit from the
bridge's extra degrees of freedom. The $6 \times 6$ coupling matrix
sits between two much larger projection matrices $A$ and $B$; any
linear transformation $M$ can be absorbed into the product $BMA$ by
redefining $B' = BM$. The optimizer has no gradient signal that would
favor distributing representation across $B$, $M$, and $A$ rather
than absorbing $M$ into $B$ entirely. This is consistent with our
Experiment~1 result: frozen-bridge RhombiLoRA matches standard LoRA
loss to within $\Delta < 0.0002$, and learnable-bridge variants do
not improve upon it. Doubling rank from 24 to 48 likewise provides
no benefit---the bottleneck is training signal, not adapter capacity.

\subsection{Why the Bridge Captures Diagnostic Structure}

Despite contributing nothing to loss, the bridge captures
task-discriminative structure because it occupies a representational
bottleneck. All information flowing from $A$ to $B$ passes through
the $6 \times 6$ bridge matrix $M$. During training, the gradients
that shape $M$ reflect how the task's loss landscape requires the
six rank channels to interact. If a task demands broad cross-channel
mixing (as general instruction following does), $M$ develops
off-diagonal entries; if the task is more constrained (as code
generation is), $M$ remains closer to diagonal.

This is the key distinction between \emph{capacity} and
\emph{diagnostics}. The bridge does not need to improve the adapter's
representational power to be informative about what the adapter
learned. It is a 36-parameter summary of the coupling pattern that
$A$ and $B$ jointly discovered---readable without running inference,
without evaluation on held-out data, and without access to the
training set. The coupling magnitude (Fiedler value) differentiates
tasks: alpaca adapters develop $2.2\times$ more cross-channel
coupling than code adapters (Fiedler 0.040 vs.\ 0.018). The coupling
structure (eigenspectrum shape) is conserved: cosine similarity
between task eigenspectra exceeds 0.999. Tasks modulate the
\emph{intensity} of bridge coupling while preserving its
\emph{architecture}.

\subsection{The Q-Projection Finding}

Across all three experimental phases, query-projection bridges are
the most informative component. Phase~0 (single-task anatomy):
Q-proj develops 40\% more coupling than V-proj ($F = 6.022$,
$p = 0.0008$). Phase~1A (task fingerprinting): Q-proj carries the
highest between-task distance (0.162 vs.\ 0.075 for K-proj and
0.072 for V-proj). Phase~2A (adapter merging): Q-proj bridges are
the most sensitive to cross-task interpolation. This ranking is
consistent across phases (Pearson $r = 0.921$ between Phase~1A task
discrimination and Phase~2A merge sensitivity).

The mechanism is architectural. The query projection computes
\emph{what to attend to}---the most task-specific computation in
the attention head. Key projections compute indexing; value
projections carry content. Queries encode the task's question about
input structure, which varies maximally across task types. The bridge
at the Q-projection captures this variation because it sits at the
point of maximum task-dependent divergence in the information flow.

An unexpected secondary finding: output projections surpass query
projections in the final five transformer layers (Layer~24 O-proj
achieves the single highest between-task distance of 0.3188 across
all 112 adapter positions). Late-layer output projections reformat
attention results for downstream prediction, and this reformatting
becomes strongly task-specific as the model approaches its output.
The depth gradient is significant (Spearman $\rho = 0.701$,
$p < 0.0001$), with a $10.3\times$ dynamic range between the most
and least discriminative bridge positions.

\subsection{Practical Implications}

A leave-one-out SVM trained on the 28 Q-projection bridge matrices
of a single adapter classifies task type at 84.5\% accuracy across
three tasks (chance: 33.3\%). This fingerprint costs 1{,}008
parameters ($28 \times 36$)---smaller than a single hidden state
vector in Qwen2.5-7B (4{,}096 dimensions). Including the 28
O-projection bridges (2{,}016 parameters total) yields 83.3\%
accuracy with improved performance on the hardest class (math:
66.1\% vs.\ 50.0\% with Q-proj alone).

For adapter composition, bridge-level interpolation preserves
eigenspectrum structure at all mixing ratios (cosine $> 0.999$),
but the Fiedler trajectory is non-linear for structurally
dissimilar task pairs ($R^2 = 0.735$ for alpaca$\leftrightarrow$math
vs.\ $R^2 = 0.956$ for alpaca$\leftrightarrow$code). The
code$\leftrightarrow$math interpolation exhibits a Fiedler dip at
$\alpha \approx 0.2$, consistent with destructive interference
between bridge patterns that differ in structure rather than
magnitude alone. At 10\% cross-task contamination ($\alpha = 0.1$),
all pairs preserve the dominant task's structure with cosine 1.0000.
This provides an empirical safe mixing ratio---conservative, but
grounded.

These results suggest a practical architecture for adapter management:
store only the 28 Q-projection bridges per task (1{,}008 parameters)
as a task fingerprint; use these for routing incoming requests to the
correct adapter without inference. K-proj and V-proj bridges are
structurally conserved across tasks and can be shared. Q-proj bridges
require task-specific handling during composition.

\subsection{Limitations}
\label{sec:limitations}

\textbf{Single model family.} All experiments use Qwen2.5-7B-Instruct.
The Q-projection dominance and depth gradient may reflect this
architecture's attention structure rather than a universal property
of transformer adapters. Extending to LLaMA, Mistral, or non-decoder
architectures is necessary before claiming generality.

\textbf{Three tasks.} Three-way classification at 84.5\% is
encouraging but limited. With only three task categories (general,
code, math), the decision boundary is low-dimensional. Scaling to
10--50 tasks would test whether bridge fingerprints remain
discriminative as the classification problem grows harder, or whether
structurally similar tasks (e.g., code generation vs.\ code
debugging) produce indistinguishable bridges.

\textbf{Fixed bridge rank.} The bridge is $6 \times 6$ in all
experiments, motivated by the rhombic dodecahedron's six face pairs.
Higher bridge ranks ($8 \times 8$, $12 \times 12$) might capture
finer-grained coupling structure at the cost of increased parameter
count. Lower ranks ($3 \times 3$, the cubic baseline) show
$4.6\times$ less coupling in Experiment~1, but whether they retain
fingerprinting ability is untested.

\textbf{Unequal training durations.} Alpaca trained for 10K steps,
code for 1{,}500, and math for 2{,}000. The monotonic relationship
between training duration and bridge distance from initialization
(alpaca: 0.199, math: 0.069, code: 0.032) means that some of the
between-task separation may reflect training duration rather than
task structure. The step-0 control (all distances exactly 0.000000)
confirms that fingerprints emerge from training, but a fully
controlled comparison would train all tasks for the same number of
steps.

\textbf{No generative evaluation.} Bridge merging was evaluated
structurally (eigenspectrum cosine, Fiedler trajectories) but not
generatively. Whether swapping bridges between adapters produces
usable output---whether the 36-parameter bridge carries enough
task information to redirect a fixed $A/B$ projection---remains
unvalidated. This is the critical missing experiment for the adapter
composition claim.

\textbf{Classification methodology.} Leave-one-out SVM on 84
samples (28 per task) is a small-sample evaluation. The classifier
uses flattened $6 \times 6 = 36$-dimensional feature vectors without
domain-informed feature engineering. More sophisticated approaches
(spectral features, graph-theoretic metrics, GL($r$)-equivariant
representations~\cite{putterman2024learning}) might improve
accuracy but would complicate the interpretability argument.

\subsection{Connection to Lattice Topology}
\label{sec:lattice-connection}

The bridge is where the lattice geometry of
Papers~1--2~\cite{bielec2026rhombic, bielec2026weighted} meets
neural architecture. The rhombic dodecahedron's six face pairs
define six channel pairs in the bridge matrix; FCC lattice
connectivity determines the initialization topology. Paper~1
established that FCC lattices achieve $2.4\times$ algebraic
connectivity over cubic lattices at matched node count. Paper~2
showed this advantage amplifies to $6.1\times$ under structured
edge weights. The bridge is the site where these lattice properties
enter the neural network.

In practice, the geometric motivation provides a principled bridge
dimension (6, from the RD's face-pair count) and initialization
structure (FCC-topology coupling), but training erases the
geometric semantics. The bridge does not learn to associate
specific channels with specific rhombic dodecahedron directions
(L-001 in the experimental record: co-planar/cross-planar ratio
$= 1.002$, $p = 0.474$). What survives is the algebraic
connectivity advantage: six channels with FCC-topology
initialization develop more cross-channel coupling than three
channels with cubic initialization ($4.6\times$ at 1.5B scale,
$1.73\times$ at 7B). The topology determines \emph{how much}
coupling the bridge develops; the task determines \emph{what for}.

This separation---geometric structure providing the capacity,
training providing the content---mirrors the lattice findings at
larger scale. The FCC lattice does not prefer specific directions;
it provides uniformly higher connectivity that any application can
exploit. The bridge inherits this property.

% ===================================================================
\section{Conclusion}
\label{sec:conclusion}

We introduced a learnable $6 \times 6$ coupling matrix---the
\emph{bridge}---between the $A$ and $B$ projections of a LoRA
adapter, adding 36 parameters per layer. The bridge does not improve
fine-tuning loss: at rank 24, its parameters are a negligible
fraction of total adapter capacity, and the optimization landscape
does not benefit from the additional degrees of freedom. This is
not a better adapter. It is a new lens on adapter behavior.

The bridge reveals what LoRA learns, in 36 parameters per layer.
Tasks that require diverse cross-channel mixing (general instruction
following) produce bridges with $2.2\times$ higher algebraic
connectivity than tasks with constrained structure (code generation).
A leave-one-out SVM trained on 28 query-projection bridges---1{,}008
parameters total---classifies task type at 84.5\% accuracy without
inference or held-out evaluation. The fingerprint emerges entirely
from training: at step~0, all bridges are identical; by step~1{,}500,
binary classification between the two most similar tasks reaches 93\%.
The query projection is the optimal fingerprint location because it
encodes the most task-specific aspect of attention: what to look for.

Bridge-level adapter composition preserves eigenspectrum structure
across all mixing ratios (cosine $> 0.999$), but non-linear Fiedler
trajectories in structurally dissimilar task pairs reveal destructive
interference invisible to full-weight merging diagnostics. The bridge
decomposes adapter behavior into a \emph{projection} component
($A$ and $B$, millions of parameters) and a \emph{coupling} component
($M$, 36 parameters)---and the coupling component alone carries
enough structure to identify the task, detect interference during
merging, and differentiate attention head components by function.

The present study is limited to one model family, three tasks, and
structural evaluation without generative validation. Scaling to more
tasks will test whether bridge fingerprints remain discriminative in
higher-dimensional classification problems. Higher bridge ranks may
capture finer coupling structure. The most immediate practical
direction is bridge-based adapter routing: given a library of
task-specific adapters, select the correct one by comparing a query
adapter's 1{,}008-parameter Q-projection fingerprint against stored
fingerprints, without running any of the adapters at inference time.

\paragraph{Availability.}
Source code: \url{https://github.com/promptcrafted/rhombic}.\\
Package: \url{https://pypi.org/project/rhombic/}.\\
Bridge matrices and analysis scripts: included in the repository
under \texttt{results/} and \texttt{scripts/}.
